\documentclass[11pt]{article}
\usepackage[margin=1in]{geometry}
\usepackage[T1]{fontenc}
\usepackage{graphicx}
\usepackage{hyperref}
\usepackage{enumitem}
\setlength{\parskip}{0.5em}
\setlength{\parindent}{0pt}

\title{graph neural networks for drug discovery}
\author{arXiv GraphRAG --- phase-6 deep research}
\date{\today}

\begin{document}
\maketitle

\begin{abstract}
The review is highly coherent and well-grounded but misses the recent paradigm shift toward molecular foundation models and the integration of LLMs with graph representations.
\end{abstract}

\section{From Topological Representations to Generative Intelligence: A Review of Graph Neural Networks in Drug Discovery}

The discovery of novel therapeutic agents remains one of the most computationally expensive and time-consuming endeavors in modern medicine, often stymied by the vastness of chemical space and the complexity of molecular interactions. Traditional computational methods frequently struggle to capture the nuanced, non-Euclidean relationships inherent in chemical structures, where the arrangement of atoms and bonds dictates biological function. Graph Neural Networks (GNNs) have emerged as a transformative solution to this problem, providing a mathematical framework that treats molecules as natural graphs. By learning directly from these topological representations, GNNs allow for a more sophisticated understanding of molecular properties, offering the potential to accelerate lead optimization and reduce the reliance on costly experimental synthesis.

This review traces the historical and technical trajectory of GNNs in the pharmaceutical sciences, beginning with the foundational principles of graph learning and molecular representation. We first examine the early advances that revolutionized molecular property prediction and the application of neural networks to force fields, establishing the baseline for predictive modeling. The discussion then moves toward the shift from predictive to creative tasks, focusing on generative models for \textit{de novo} molecular design. To address the physical realities of molecular biology, we explore the integration of geometric, 3D, and equivariant graph neural networks that account for spatial orientation. The review further investigates advanced architectures, specifically the adoption of attention mechanisms, transformers, and self-supervision techniques to enhance feature extraction. We conclude by addressing the critical requirements for interpretability, explainability, and structural learning, before surveying the current frontier of large-scale trends and emerging directions that are shaping the future of AI-driven drug discovery.

\subsection{Foundations of Graph Learning and Molecular Representation}

The evolution of graph learning has transitioned from spectral-based methodologies to spatial convolutional architectures. Early foundational work focused on spectral networks and locally connected networks on graphs (\texttt{arXiv:1312.6203}), providing a basis that has been widely extended and compared in subsequent literature (\texttt{arXiv:2211.03666}). This spectral approach was later contrasted with convolutional neural networks designed to operate directly on graphs for the purpose of predicting molecular fingerprints (\texttt{arXiv:1509.09292}).

A pivotal shift toward more scalable spatial methods was marked by the introduction of Graph Convolutional Networks (GCNs) for semi-supervised classification (\texttt{arXiv:1609.02907}). GCNs offered a framework that could be compared against both spectral models (\texttt{arXiv:1312.6203}) and other convolutional approaches (\texttt{arXiv:1606.09375}). This lineage of convolutional learning diverged into two primary architectural improvements: inductive learning and attention mechanisms. GraphSAGE addressed the need for inductive representation learning on large graphs by utilizing sampling and aggregation techniques, which proved effective for node classification tasks such as identifying protein functions (\texttt{arXiv:1706.02216}). Simultaneously, Graph Attention Networks (GATs) extended the GCN architecture (\texttt{arXiv:1609.02907}) by implementing masked self-attentional layers (\texttt{arXiv:1710.10903}), allowing for more sophisticated feature weighting.

Despite these advancements, the structural capacity of these models remains a subject of critical scrutiny. Research into the discriminative power of GNNs has demonstrated that the ability of architectures like GCNs and GraphSAGE to perform graph classification is bounded; specifically, their effectiveness can be measured against an architecture equivalent to the Weisfeiler-Lehman (WL) test (\texttt{arXiv:1810.00826}). This critique highlights a fundamental limitation in how standard GNNs perceive graph topology and distinguish complex structures. To facilitate the practical implementation of these diverse architectures—including GCNs (\texttt{arXiv:1609.02907}), GATs (\texttt{arXiv:1710.10903}), and GraphSAGE (\texttt{arXiv:1706.02216})—the development of optimized frameworks like PyTorch Geometric has become essential for fast graph representation learning (\texttt{arXiv:1903.02428}).

\subsection{Early Advances in Molecular Property Prediction and Force Fields}

Early research in molecular machine learning focused on transitioning from traditional fingerprints to representations capable of capturing complex quantum-mechanical and structural properties. A foundational challenge in this domain was the modeling of atomization energies using nuclear charges and atomic positions, with early efforts achieving a mean absolute error of approximately 10 kcal/mol [\texttt{arXiv:1109.2618}]. While these early models established the feasibility of using machine learning for energy prediction, subsequent developments significantly improved both accuracy and the depth of chemical insight.

A major shift occurred with the introduction of deep tensor neural networks (DTNN), which moved beyond simple energy prediction to model quantum-mechanical observables, including local chemical potentials and isomer energies [\texttt{arXiv:1609.08259}]. The DTNN approach demonstrated high precision, achieving an accuracy of 1 kcal/mol, and provided a distinct advantage through spatially and chemically resolved explainability [\texttt{arXiv:1609.08259}]. This emphasis on physical accuracy is further reflected in the construction of energy-conserving molecular force fields. By utilizing conformational geometries to predict potential energy surfaces, researchers achieved high-fidelity results, specifically 0.3 kcal mol⁻¹ for energies and 1 kcal mol⁻¹ Å⁻¹ for atomic forces [\texttt{arXiv:1611.04678}].

Parallel to these quantum-mechanical advancements, the field expanded into biological applications through multitask and structure-based learning. Rather than focusing on single properties, multitask networks enabled the prediction of measurements across more than 200 biological targets [\texttt{arXiv:1502.02072}]. This capability was complemented by AtomNet, which utilized 3D structural target information to predict small molecule bioactivity, achieving an AUC greater than 0.9 on 57.8\% of targets in the DUDE benchmark [\texttt{arXiv:1510.02855}].

A critical evolution in representation was the move toward molecular graph convolutions, which allowed models to move "beyond fingerprints" by treating molecules as undirected graphs comprising atoms, bonds, and distances [\texttt{arXiv:1603.00856}]. This graph-based approach provided a more flexible framework for ligand-based virtual screening compared to earlier fixed-descriptor methods. To standardize these diverse tasks—ranging from quantum mechanical to biophysical properties—the MoleculeNet benchmark was established [\texttt{arXiv:1703.00564}]. MoleculeNet provided a unified framework using learnable representations and physics-aware featurizations, which has since become a cornerstone for evaluating the progress of molecular machine learning [\texttt{arXiv:1703.00564}]. Collectively, these early works transitioned the field from basic energy estimation to sophisticated, graph-based, and multitask architectures capable of navigating both the physical and biological complexities of drug discovery.

\subsection{Generative Models for De Novo Molecular Design}

The evolution of de novo molecular design has transitioned from sequence-based representations to increasingly sophisticated graph-structured generative frameworks. Early approaches focused on converting discrete molecular representations into multidimensional continuous latent spaces to facilitate optimization of chemical properties [\texttt{arXiv:1610.02415}]. This concept was expanded through sequence-based models, such as those using SMILES strings [\texttt{arXiv:1705.10843}] or context-free grammar parse trees [\texttt{arXiv:1703.01925}], and reinforcement learning strategies designed to target specific biological activities, such as dopamine receptor type 2 activity [\texttt{arXiv:1704.07555}].

However, a significant shift in the literature reflects a move toward direct graph generation to better capture molecular topology. Early graph-based variational autoencoders (VAEs), such as GraphVAE, utilized probabilistic fully-connected graphs [\texttt{arXiv:1802.03480}], while other models focused on learning the general structure and attributes of synthetic molecular graphs through graph neural networks [\texttt{arXiv:1803.03323}]. To address the complexities of chemical validity, MolGAN introduced an implicit generative model using Generative Adversarial Networks (GANs) adapted for 2D molecular graphs, achieving high validity rates [\texttt{arXiv:1805.11973}]. This approach builds upon relational graph convolutional networks [\texttt{arXiv:1703.06103}] to handle the relational nature of molecular data.

A critical tension exists between the simplicity of sequence models and the structural integrity of graph models. While sequence-based models like ORGAN [\texttt{arXiv:1705.10843}] allow for arbitrary metric optimization, they face scrutiny regarding their representational limits [\texttt{arXiv:1802.04364}]. In contrast, more recent architectures prioritize chemical constraints and hierarchical structures. For instance, the Junction Tree VAE (JT-VAE) utilizes a graph message passing network to generate molecules via a tree-structured scaffold of chemical substructures, providing a more robust method for molecular graph generation and optimization [\texttt{arXiv:1802.04364}].

Furthermore, the field has moved toward goal-directed generation. The Graph Convolutional Policy Network (GCPN) exemplifies this by optimizing for drug-likeness, synthetic accessibility, and chemical valency, reporting a 184\% improvement in constrained property optimization [\texttt{arXiv:1806.02473}]. Other constrained graph VAEs focus on ensuring that generated molecules match the semantic statistics of original datasets [\texttt{arXiv:1805.09076}]. Collectively, these advancements demonstrate a trajectory from continuous latent space optimization [\texttt{arXiv:1610.02415}] toward highly constrained, goal-oriented graph architectures that balance chemical validity with specific pharmacological objectives [\texttt{arXiv:1802.04364}, \texttt{arXiv:1806.02473}].

\subsection{Geometric, 3D, and Equivariant Graph Neural Networks}

The evolution of graph neural networks (GNNs) for molecular modeling has transitioned from simple topological representations to sophisticated architectures that explicitly incorporate 3D geometry and physical symmetries. Early advancements in this domain focused on modeling quantum interactions through continuous-filter convolutional neural networks. SchNet [\texttt{arXiv:1706.08566}, \texttt{arXiv:1712.06113}] introduced a framework utilizing continuous-filter convolutional layers to process 3D atomic positions, enabling the prediction of total energy, interatomic forces, and potential-energy surfaces. While effective for modeling quantum-mechanical properties, these architectures primarily rely on atomic distances rather than full geometric orientation.

To address the limitations of distance-only models, subsequent research integrated directional and rotational information. DimeNet [\texttt{arXiv:2003.03123}] introduced directional message passing, utilizing both atom distances and spatial directions within 3D graphs to predict quantum mechanical properties. This approach moves beyond the scalar-based interactions of SchNet by incorporating the coordinate space explicitly. A more rigorous treatment of symmetry is found in Tensor Field Networks [\texttt{arXiv:1802.08219}], which achieve local equivariance to 3D rotations, translations, and permutations when processing 3D point clouds. Similarly, Cormorant [\texttt{arXiv:1906.04015}] employs a rotationally covariant architecture using tensor products and Clebsch-Gordan decomposition to learn atomic potential energy surfaces and ground state properties from conformational geometries.

The field has recently converged on the necessity of formal equivariance to ensure physical consistency. E(n)-Equivariant Graph Neural Networks (EGNNs) [\texttt{arXiv:2102.09844}] provide a generalized framework that is equivariant to rotations, translations, reflections, and permutations, facilitating robust representation learning and dynamical systems modeling. This principle of equivariance is further refined in Equivariant Graph Attention Networks [\texttt{arXiv:2202.09891}], which apply attention mechanisms to 3D Cartesian coordinates to predict quantum mechanical properties and macromolecular structures.

The utility of these geometric representations extends from property prediction to molecular design. For instance, the Ligand Neural Network (L-Net) [\texttt{arXiv:2104.08474}] leverages 3D structures—including topological data and hydrogen atom coordinates—to drive deep generative models. This allows for the generation of drug-like molecules with specific targets, such as ABL1 kinase inhibition. Collectively, these developments demonstrate a clear trajectory: from SchNet’s continuous-filter convolutions [\texttt{arXiv:1706.08566}] to the highly symmetric, equivariant architectures of EGNNs [\texttt{arXiv:2102.09844}] and Cormorant [\texttt{arXiv:1906.04015}], the field has moved toward architectures that fundamentally respect the geometric and physical constraints of 3D molecular space.

\subsection{Advanced Architectures: Attention, Transformers, and Self-Supervision}

Recent advancements in graph neural networks (GNNs) for drug discovery have shifted from simple message passing toward sophisticated architectures that leverage hierarchical structures, self-supervision, and geometric information. A foundational step in this evolution is the development of hierarchical representation learning. \texttt{arXiv:1806.08804} introduced differentiable pooling to capture graph-level hierarchies, a method primarily evaluated through graph classification accuracy. This hierarchical approach is further refined in the context of molecular property prediction (MPP) by \texttt{arXiv:2208.13994}, which proposes HiGNN. Unlike the purely structural focus of \texttt{arXiv:1806.08804}, HiGNN employs feature-wise attention and a co-representation of molecular graphs alongside chemically synthesizable BRICS fragments, enhancing explainability through a molecule-fragment similarity mechanism to predict druggability and bioactivities.

The integration of self-supervision has become a critical strategy for overcoming data scarcity in molecular datasets. \texttt{arXiv:1905.12265} established comprehensive pre-training strategies for both individual nodes and entire graphs, demonstrating efficacy in downstream molecular property and protein function prediction tasks using ROC-AUC metrics. This paradigm was significantly extended by \texttt{arXiv:2007.02835}, which introduced a Self-Supervised Graph Transformer. By integrating message passing networks into a Transformer-style architecture, this model achieves superior molecular representation learning on 2D graphs across 11 benchmarks, moving beyond the standard GNN frameworks used in earlier pre-training studies.

While many models focus on 2D topological structures, recent work has sought to incorporate higher-dimensional spatial data. \texttt{arXiv:2109.11730} addresses this by utilizing a dual-view geometric message passing network (GeomMPNN) for contrastive learning. By representing molecules as both 2D and 3D graphs, this approach captures essential geometric information that purely topological models may overlook. This emphasis on structural nuance is also reflected in \texttt{arXiv:2312.16855}, which enhances MPP by incorporating a molecular similarity graph (MSG) constructed from molecular fingerprints, effectively augmenting the standard molecular graph with learned structural relationships.

Finally, the scope of advanced architectures has expanded from prediction to generative tasks. \texttt{arXiv:2304.05498} introduces GraphGANFed, a federated generative framework using Graph Convolutional Networks (GCNs) to generate molecular graphs. While previous works like \texttt{arXiv:2007.02835} and \texttt{arXiv:2109.11730} focus on the discriminative task of property prediction, GraphGANFed prioritizes generative metrics, achieving a novelty of 100 and diversity greater than 0.9. Collectively, these developments demonstrate a trajectory from basic graph classification toward multi-modal, self-supervised, and generative models that integrate chemical intuition—such as BRICS fragments or geometric 3D data—into the deep learning pipeline.

\subsection{Interpretability, Explainability, and Structural Learning}

The advancement of Graph Neural Networks (GNNs) in drug discovery is increasingly defined by the transition from black-box predictive models to systems capable of providing structural rationale for molecular properties. Current research diverges between post-hoc interpretability methods and the integration of hierarchical structural information directly into the model architecture.

A primary focus in the literature is the quantitative evaluation of explainability methods applied to standard architectures like GraphNet and CMPNN. For instance, \texttt{arXiv:2107.04119} provides a quantitative assessment of interpretability for molecular property prediction, specifically comparing GradInput and Integrated Gradients (IG). While these methods attempt to assign importance to specific graph components, they often operate as post-hoc explanations rather than intrinsic structural learners. In contrast, more recent approaches attempt to bake interpretability into the representation itself. \texttt{arXiv:2304.06253} enhances compound property and activity prediction by utilizing multiple molecular graph representations, moving beyond simple atom-level graphs to incorporate reduced molecular graph representations that capture higher-level substructures. This approach allows for an interpretation analysis of both individual features and potential substructures, bridging the gap between local atomistic detail and global molecular topology.

This trend toward structural hierarchy is further refined in the development of FragNet, which introduces a multi-scale interpretability framework. Unlike models that focus on a single level of granularity, FragNet (\texttt{arXiv:2410.12156}) provides four distinct levels of interpretability: atoms, bonds, molecular fragments, and the connections between those fragments. This hierarchical depth offers a more granular view of molecular property prediction than the dual-representation approach seen in \texttt{arXiv:2304.06253}.

Another significant methodology involves using sparsity to enforce explainability. For the specific task of drug-protein binding prediction, \texttt{arXiv:2309.12906} introduces SLGNN (Sparse Learning to Graph Neural Networks). By employing a chemical-substructure-based graph where nodes represent entire substructures rather than individual atoms, and applying a generalized fused lasso, the model identifies connected subgraphs through sparse learning. This method contrasts with the gradient-based approaches in \texttt{arXiv:2107.04119} by using structural sparsity to inherently select the most relevant chemical motifs, rather than merely calculating importance scores after training.

Finally, the field is addressing the fundamental design of graph-structured data to ensure models capture invariant and equivariant molecular characteristics. \texttt{arXiv:2411.12156} explores molecular graph generation by utilizing GNNs with advanced high-order pooling functions, incorporating both atom-bond structures and precise atom positions. While the aforementioned papers focus on interpreting existing predictions, this work emphasizes the importance of the underlying representation's geometric accuracy to ensure that the learned features are physically meaningful. Together, these works demonstrate a shift from simple attribution methods toward complex, multi-level structural learning that mirrors the hierarchical nature of chemical entities.

\subsection{The Current Frontier: Large-Scale Surveys and Emerging Trends}

The landscape of GNNs in drug discovery is characterized by a transition from basic molecular property prediction toward complex, multi-scale modeling and generative frameworks. Recent surveys provide a foundational overview of this evolution. \texttt{arXiv:2506.01302} and \texttt{arXiv:2509.07887} offer broad taxonomies covering diverse tasks, ranging from retrosynthesis and drug repositioning to microbiome interaction prediction. Notably, \texttt{arXiv:2506.06915} advances this field by providing a model-centric taxonomy that bridges mathematical GNN theory—specifically spectral versus spatial methods—with physical inductive biases such as invariance and equivariance. This work is particularly strong in its ability to map classic architectures, such as ChebNet and GCN, into the modern Message Passing Neural Network (MPNN) framework. However, a significant gap remains in the literature; while \texttt{arXiv:2506.06915} establishes these theoretical connections, it lacks a critical evaluation of the trade-offs between the geometric complexity of advanced models (such as 3rd-order interactions in GemNet) and the computational scalability necessary for large-scale, high-throughput drug screening.

This tension between architectural complexity and utility is also evident in the debate over predictive modeling. While \texttt{arXiv:2411.13688} investigates the efficacy of Graph Isomorphism Networks (GINs) by comparing them against classical feature-extraction techniques, such as extended-connectivity fingerprints (ECFPs) and physicochemical-descriptor vectors (PDVs) for activity-cliff prediction, other researchers are seeking to enhance GNNs through architectural fusion. For instance, \texttt{arXiv:2507.03430} moves beyond simple graph representations by employing a Multi-Level Fusion Graph Neural Network (MLFGNN) that integrates Graph Attention Networks (GAT) and Graph Transformers with 2D graphs and molecular fingerprints to improve classification and regression accuracy.

The frontier also extends into generative and relational modeling. Rather than treating prediction and generation as distinct stages, \texttt{arXiv:2406.03278} explores a novel paradigm by using GNN property predictors—specifically those targeting electronic properties and logP—directly as molecule generators. This shift toward active design is complemented by efforts to model complex biological networks. While much of the literature focuses on single-molecule properties, \texttt{arXiv:2502.09335} employs a Graph Diffusion Network (GDNDGP) to predict drug-gene associations, leveraging meta-path-based homogeneous graphs to represent drug-drug and gene-gene relationships. Together, these developments indicate a move toward models that are not only more mathematically rigorous but also more deeply integrated into the broader biological context of drug-target and drug-gene interactions.

\subsection{Tensions and Open Problems}

A primary tension in molecular design lies in the architectural philosophy of generative modeling: whether to optimize for local atomistic connectivity or global structural validity. Early reinforcement learning approaches, such as the Objective-Reinforced Generative Adversarial Networks described in \texttt{arXiv:1705.10843}, attempt to guide generation through reward-based feedback on specific properties. However, this method often struggles with the fundamental validity of the resulting chemical structures. In contrast, the Junction Tree Variational Autoencoder proposed in \texttt{arXiv:1802.04364} shifts the focus toward a hierarchical approach, generating molecules by assembling chemically valid sub-structures (junction trees) rather than individual atoms, thereby ensuring structural integrity. Recent investigations in \texttt{arXiv:2406.03278} further complicate this dichotomy by suggesting that the gap between prediction and generation can be bridged by repurposing high-performing GNN property predictors as generative engines, potentially bypassing the need for complex, explicit generative architectures altogether.

Another significant debate concerns the optimal level of geometric abstraction required for accurate property prediction. There is a clear divide between models that prioritize 2D topological features and those that incorporate 3D spatial configurations. Traditional methods, such as the convolutional networks on graphs explored in \texttt{arXiv:1509.09292} and the message-passing frameworks in \texttt{arXiv:1704.01212}, focus heavily on the connectivity and local chemical environments of atoms. While effective for many tasks, these models often lack the ability to capture the nuances of three-dimensional conformations. This limitation is highlighted in \texttt{arXiv:2305.13315}, which contrasts the topological focus of earlier works with the continuous-filter convolutions of SchNet (\texttt{arXiv:1706.08566}) and the directional message passing of \texttt{arXiv:2003.03123}, both of which attempt to integrate 3D geometric information. The tension remains whether the computational overhead of modeling full 3D coordinates is justified, or if sophisticated 2D graph representations can implicitly capture the necessary structural information.

Finally, the field faces a fundamental disagreement regarding the representational capacity and design robustness of GNNs. While foundational works like the Graph Convolutional Networks in \texttt{arXiv:1609.02907} and Graph Attention Networks in \texttt{arXiv:1710.10903} established the standard for neighborhood aggregation, recent critiques suggest these designs may be insufficient for complex graph structures. For instance, \texttt{arXiv:2411.07269} argues that many standard GNN designs suffer from inherent flaws in how they aggregate information, potentially limiting their ability to capture higher-order structural motifs. This is echoed by the theoretical inquiries in \texttt{arXiv:1810.00826}, which questions the expressive power of these architectures compared to the Weisfeiler-Lehman test. Consequently, a tension exists between the widespread adoption of "standard" message-passing layers and the emerging need for more sophisticated, hierarchical, or structurally aware architectures that can move beyond simple local aggregations.

\subsubsection{Outlook}

The frontier of molecular graph learning is rapidly shifting from simple property prediction toward a holistic integration of multi-scale geometric reasoning and generative autonomy. As highlighted by the most recent developments in 2025, the field is moving toward multi-level fusion architectures (\texttt{arXiv:2507.03430}) and geometric graph-based models for complex tasks like drug-target affinity prediction (\texttt{arXiv:2509.13476}). The next generation of models will likely transcend the "atom-by-atom" paradigm, instead focusing on the seamless fusion of quantum-mechanical insights, 3D conformational dynamics, and structural similarity-driven representations to achieve truly predictive and generative intelligence in drug discovery.

\section*{References}
\begin{enumerate}\small
  \item Fast and Accurate Modeling of Molecular Atomization Energies with
  Machine Learning (2011). \texttt{arXiv:1109.2618}
  \item Spectral Networks and Locally Connected Networks on Graphs (2013). \texttt{arXiv:1312.6203}
  \item Auto-Encoding Variational Bayes (2013). \texttt{arXiv:1312.6114}
  \item Adam: A Method for Stochastic Optimization (2014). \texttt{arXiv:1412.6980}
  \item Learning Phrase Representations using RNN Encoder-Decoder for
  Statistical Machine Translation (2014). \texttt{arXiv:1406.1078}
  \item Convolutional Networks on Graphs for Learning Molecular Fingerprints (2015). \texttt{arXiv:1509.09292}
  \item Gated Graph Sequence Neural Networks (2015). \texttt{arXiv:1511.05493}
  \item Massively Multitask Networks for Drug Discovery (2015). \texttt{arXiv:1502.02072}
  \item Deep Residual Learning for Image Recognition (2015). \texttt{arXiv:1512.03385}
  \item AtomNet: A Deep Convolutional Neural Network for Bioactivity Prediction
  in Structure-based Drug Discovery (2015). \texttt{arXiv:1510.02855}
  \item Automatic chemical design using a data-driven continuous representation
  of molecules (2016). \texttt{arXiv:1610.02415}
  \item Semi-Supervised Classification with Graph Convolutional Networks (2016). \texttt{arXiv:1609.02907}
  \item Molecular Graph Convolutions: Moving Beyond Fingerprints (2016). \texttt{arXiv:1603.00856}
  \item Quantum-Chemical Insights from Deep Tensor Neural Networks (2016). \texttt{arXiv:1609.08259}
  \item Convolutional Neural Networks on Graphs with Fast Localized Spectral
  Filtering (2016). \texttt{arXiv:1606.09375}
  \item Machine Learning of Accurate Energy-Conserving Molecular Force Fields (2016). \texttt{arXiv:1611.04678}
  \item Geometric deep learning: going beyond Euclidean data (2016). \texttt{arXiv:1611.08097}
  \item Neural Message Passing for Quantum Chemistry (2017). \texttt{arXiv:1704.01212}
  \item MoleculeNet: A Benchmark for Molecular Machine Learning (2017). \texttt{arXiv:1703.00564}
  \item Grammar Variational Autoencoder (2017). \texttt{arXiv:1703.01925}
  \item Attention Is All You Need (2017). \texttt{arXiv:1706.03762}
  \item SchNet: A continuous-filter convolutional neural network for modeling
  quantum interactions (2017). \texttt{arXiv:1706.08566}
  \item Objective-Reinforced Generative Adversarial Networks (ORGAN) for
  Sequence Generation Models (2017). \texttt{arXiv:1705.10843}
  \item Graph Attention Networks (2017). \texttt{arXiv:1710.10903}
  \item Inductive Representation Learning on Large Graphs (2017). \texttt{arXiv:1706.02216}
  \item Molecular De Novo Design through Deep Reinforcement Learning (2017). \texttt{arXiv:1704.07555}
  \item In silico generation of novel, drug-like chemical matter using the LSTM
  neural network (2017). \texttt{arXiv:1712.07449}
  \item Modeling Relational Data with Graph Convolutional Networks (2017). \texttt{arXiv:1703.06103}
  \item SchNet - a deep learning architecture for molecules and materials (2017). \texttt{arXiv:1712.06113}
  \item Junction Tree Variational Autoencoder for Molecular Graph Generation (2018). \texttt{arXiv:1802.04364}
  \item How Powerful are Graph Neural Networks? (2018). \texttt{arXiv:1810.00826}
  \item GraphVAE: Towards Generation of Small Graphs Using Variational
  Autoencoders (2018). \texttt{arXiv:1802.03480}
  \item BERT: Pre-training of Deep Bidirectional Transformers for Language
  Understanding (2018). \texttt{arXiv:1810.04805}
  \item Graph Convolutional Policy Network for Goal-Directed Molecular Graph
  Generation (2018). \texttt{arXiv:1806.02473}
  \item Tensor field networks: Rotation- and translation-equivariant neural
  networks for 3D point clouds (2018). \texttt{arXiv:1802.08219}
  \item Learning Deep Generative Models of Graphs (2018). \texttt{arXiv:1803.03324}
  \item MolGAN: An implicit generative model for small molecular graphs (2018). \texttt{arXiv:1805.11973}
  \item Hierarchical Graph Representation Learning with Differentiable Pooling (2018). \texttt{arXiv:1806.08804}
  \item Constrained Graph Variational Autoencoders for Molecule Design (2018). \texttt{arXiv:1805.09076}
  \item Syntax-Directed Variational Autoencoder for Structured Data (2018). \texttt{arXiv:1802.08786}
  \item Analyzing Learned Molecular Representations for Property Prediction (2019). \texttt{arXiv:1904.01561}
  \item Strategies for Pre-training Graph Neural Networks (2019). \texttt{arXiv:1905.12265}
  \item Fast Graph Representation Learning with PyTorch Geometric (2019). \texttt{arXiv:1903.02428}
  \item Cormorant: Covariant Molecular Neural Networks (2019). \texttt{arXiv:1906.04015}
  \item Molecular Property Prediction: A Multilevel Quantum Interactions
  Modeling Perspective (2019). \texttt{arXiv:1906.11081}
  \item Molecular graph generation with Graph Neural Networks (2020). \texttt{arXiv:2012.07397}
  \item Graph Neural Network Architecture Search for Molecular Property
  Prediction (2020). \texttt{arXiv:2008.12187}
  \item Utilising Graph Machine Learning within Drug Discovery and Development (2020). \texttt{arXiv:2012.05716}
  \item Learn molecular representations from large-scale unlabeled molecules for
  drug discovery (2020). \texttt{arXiv:2012.11175}
  \item Multi-View Graph Neural Networks for Molecular Property Prediction (2020). \texttt{arXiv:2005.13607}
  \item Directional Message Passing for Molecular Graphs (2020). \texttt{arXiv:2003.03123}
  \item Self-Supervised Graph Transformer on Large-Scale Molecular Data (2020). \texttt{arXiv:2007.02835}
  \item Molecular Graph Generation via Geometric Scattering (2021). \texttt{arXiv:2110.06241}
  \item Quantitative Evaluation of Explainable Graph Neural Networks for
  Molecular Property Prediction (2021). \texttt{arXiv:2107.04119}
  \item Learning to design drug-like molecules in three-dimensional space using
  deep generative models (2021). \texttt{arXiv:2104.08474}
  \item Structure-aware Interactive Graph Neural Networks for the Prediction of
  Protein-Ligand Binding Affinity (2021). \texttt{arXiv:2107.10670}
  \item GeomGCL: Geometric Graph Contrastive Learning for Molecular Property
  Prediction (2021). \texttt{arXiv:2109.11730}
  \item E(n) Equivariant Graph Neural Networks (2021). \texttt{arXiv:2102.09844}
  \item Graph Neural Networks for Molecules (2022). \texttt{arXiv:2209.05582}
  \item Graph neural networks for the prediction of molecular structure-property
  relationships (2022). \texttt{arXiv:2208.04852}
  \item HiGNN: Hierarchical Informative Graph Neural Networks for Molecular
  Property Prediction Equipped with Feature-Wise Attention (2022). \texttt{arXiv:2208.13994}
  \item Application of Graph Neural Networks and graph descriptors for graph
  classification (2022). \texttt{arXiv:2211.03666}
  \item Equivariant Graph Attention Networks for Molecular Property Prediction (2022). \texttt{arXiv:2202.09891}
  \item Building explainable graph neural network by sparse learning for the
  drug-protein binding prediction (2023). \texttt{arXiv:2309.12906}
  \item Molecular Property Prediction Based on Graph Structure Learning (2023). \texttt{arXiv:2312.16855}
  \item GraphGANFed: A Federated Generative Framework for Graph-Structured
  Molecules Towards Efficient Drug Discovery (2023). \texttt{arXiv:2304.05498}
  \item 3D Molecular Geometry Analysis with 2D Graphs (2023). \texttt{arXiv:2305.13315}
  \item Enhancing Model Learning and Interpretation Using Multiple Molecular
  Graph Representations for Compound Property and Activity Prediction (2023). \texttt{arXiv:2304.06253}
  \item Using GNN property predictors as molecule generators (2024). \texttt{arXiv:2406.03278}
  \item Learning From Graph-Structured Data: Addressing Design Issues and
  Exploring Practical Applications in Graph Representation Learning (2024). \texttt{arXiv:2411.07269}
  \item Molecular Graph Representation Learning via Structural Similarity
  Information (2024). \texttt{arXiv:2409.08580}
  \item FragNet: A Graph Neural Network for Molecular Property Prediction with Four Levels of Interpretability (2024). \texttt{arXiv:2410.12156}
  \item Investigating Graph Neural Networks and Classical Feature-Extraction
  Techniques in Activity-Cliff and Molecular Property Prediction (2024). \texttt{arXiv:2411.13688}
  \item On the Scalability of GNNs for Molecular Graphs (2024). \texttt{arXiv:2404.11568}
  \item Graph Neural Networks in Modern AI-aided Drug Discovery (2025). \texttt{arXiv:2506.06915}
  \item A Survey of Graph Neural Networks for Drug Discovery: Recent Developments and Challenges (2025). \texttt{arXiv:2509.07887}
  \item Recent Developments in GNNs for Drug Discovery (2025). \texttt{arXiv:2506.01302}
  \item Graph Diffusion Network for Drug-Gene Prediction (2025). \texttt{arXiv:2502.09335}
  \item Multi-Level Fusion Graph Neural Network for Molecule Property Prediction (2025). \texttt{arXiv:2507.03430}
  \item A Geometric Graph-Based Deep Learning Model for Drug-Target Affinity Prediction (2025). \texttt{arXiv:2509.13476}
\end{enumerate}

\end{document}
